% !TEX root = ../PolycopiePremiere2627.tex
%%%%%%%%%%%%%%%%%%%%%%%%%%%%%%%%%%%%%%%%%%%%%%%%%%%%%%%%%%%%%%%%%%%%%%%%%%%%%%%%
% Contenu initial repris de Louis Paternault --- http://ababsurdo.fr
%
% Publié sous licence Creative Commons Attribution-ShareAlike 4.0 International (CC BY-SA 4.0)
% http://creativecommons.org/licenses/by-sa/4.0/deed.fr
%%%%%%%%%%%%%%%%%%%%%%%%%%%%%%%%%%%%%%%%%%%%%%%%%%%%%%%%%%%%%%%%%%%%%%%%%%%%%%%%

\chapter{Fonction affine}

\Dfn{Une \emph{fonction affine} est une fonction pouvant être exprimée sous la forme $x\mapsto ax+b$, où $a$ et $b$ sont des nombres réels. Elle est définie sur \textcolor{J1}{$\mathbb R$}.}

\Dfn{La courbe d'une fonction affine $f:x\mapsto ax+b$ est une droite, dont :
\begin{itemize}
    \item $a$ est \textcolor{J1}{le coefficient directeur}.
    \item $b$ est \textcolor{J1}{l'ordonnée à l'origine}.
\end{itemize}}

\Prp{Soit $f$ une fonction affine.
\begin{itemize}
    \item Soient $A$ et $B$ deux points de sa courbe représentative. Alors le coefficient directeur est donné par la formule :
    \(\dfrac{y_B-y_A}{x_B-x_A}\).
    \item Soient deux nombres $a$ et $b$ distincts. Alors le coefficient directeur est donné par la formule :
    \(\dfrac{f(b)-f(a)}{b-a}\).
\end{itemize}}

\Exemple{
\begin{enumerate}
    \item Déterminer le coefficient directeur de la fonction affine dont la droite passe par les points $A(2;-7)$ et $B(7;0)$.
    \item Déterminer le coefficient directeur de la fonction affine $f$ telle que $f(-8)=3$ et $f(2)=0$.
\end{enumerate}}

\Prp{Soit $f:x\mapsto ax+b$ une fonction affine. Alors :
\begin{itemize}
    \item si $a<0$, alors la fonction est \textcolor{J1}{strictement décroissante} ;
    \item si $a=0$, alors la fonction est \textcolor{J1}{constante} ;
    \item si $a>0$, alors la fonction est \textcolor{J1}{strictement croissante}.
\end{itemize}}

\Dfn{Une fonction affine permet de modéliser la croissance d'un phénomène linéaire continu :
\begin{itemize}
    \item linéaire : le taux d'accroissement de la fonction est constant ;
    \item continu (par opposition à \emph{discret}) : le phénomène peut être étudié sur un intervalle (plutôt que, par exemple, sur l'ensemble des entiers).
\end{itemize}}

\Exemple{Dire si les phénomènes suivants sont linéaires ou non, continus ou discrets.
\begin{enumerate}
    \item Évolution de l'aire d'un carré en fonction de la longueur d'un de ses côtés.
    \item Évolution de la température d'une casserole d'eau sur le feu en fonction du temps.
    \item Évolution de la somme d'argent dans une tirelire dans laquelle on ajoute 10 € tous les mois.
    \item Évolution de la TVA d'un article en fonction de son prix.
\end{enumerate}}
